Eight PID tunings were compared in an implementation study on a three-state exothermic CSTR whose cooling valve is limited in amplitude and in slew rate, under a single protocol with true-temperature metrics. The contribution is an analysis of implementation, not a new anti-windup law. Proposition~\ref{prop:latch} is general for reverse-acting loops with an open-loop-stable plant; the regime map, the comparison of tunings and the sensitivity analysis refer to the reactor and valve studied. Four conclusions follow.
\begin{enumerate}
    \item \textbf{The anti-windup implementation can decide the verdict on a tuning rule, and the slew-rate trigger is what matters.} In a reverse-acting loop, sign-blind amplitude clamping, a frequent implementation error, is worse than no anti-windup: it never freezes the integrator while it winds up and freezes it when it must unwind, latching the valve at its upper stop (Proposition~\ref{prop:latch}); on this reactor it sits $16.7$~K below the setpoint and the $\mathrm{IAE}$ is $104.123$ and $102.668$ instead of $0.940$ and $1.017$~K$\cdot$min for ZN and CC. A $2\times 2$ factorial on three setpoint steps, a load profile, a second operating point and a valve deadband attributes the recovery to a rate-aware trigger; directional sign awareness is second order. Over slew limits of $20$ to $320$~L/min$^2$, with fixed or re-optimized gains, the trigger matters for ZN and CC at every limit, for AMIGO and Optimal ITAE when the valve is slow, and not for SIMC. Back-calculation with a single $T_t$ per loop chosen by a declared criterion (minimax over the three steps, or the design step alone) stays within $2\%$ of the directional rule in IAE over the steps and is up to $10\%$ better; the usual defaults $T_t = \sqrt{T_iT_d}$ and $T_t = T_i$ are up to $37\%$ and $97\%$ worse. The directional rule needs no per-loop constant but depends on $S_{\max}$ and $\Delta t$. Comparisons that report the failure of empirical rules should state their anti-windup logic (Section~\ref{subsec:checklist}).
    \item \textbf{Noise vulnerability is set by the high-frequency controller gain, not by a sensitivity peak of the plant.} On the exact linearization ZN and CC have $\|S\|_\infty = 1.57$ and $1.75$, but $|K_p|(1+N)\approx\|CS\|_\infty$ of $490$ and $559$~L/min/K against $77$ for SIMC. Lowering the derivative-filter coefficient to $N = 3$ cuts that gain by $64\%$ at no cost in nominal error; a first-order measurement filter ($\tau_m = 0.05$~min) lowers valve travel by a factor of $1.8$ to $3.3$ ($44$ to $69\%$) at $\sigma = 0.01$~K and makes valve wear discriminate between tunings, but neither keeps the empirical rules away from the kinematic travel ceiling at $\sigma = 0.1$~K.
    \item \textbf{Imposing the sensitivity bound on the exact loop buys an in-sample gain only.} Bounding $M_s$ on the identified dead-time loop gives $M_s^{\mathrm{lin}} = 1.31$ on the exact loop. Imposing $M_s^{\mathrm{lin}} \le 1.6$ instead yields an interior optimum ($M_s^{\mathrm{lin}} = 1.20$) with $11.3\%$ lower ITAE and $3.4\%$ lower IAE on the design step, at $76\%$ higher high-frequency gain, $22.6\%$ more valve travel and no gain on out-of-sample steps. Under identical $\pm 20\%$ errors of loop gain, all time constants and delay on the exact plant, no tuning has an unstable corner, but the peak of the exact-loop optimum rises to $2.06$ against $1.63$ for the FOPTD-bounded one. Optimizing the setpoint weights jointly with the gains returns the 1-DOF controller; setpoint weighting is a trade of tracking speed for valve travel, not an optimum of the tracking objective. The FOPTD-bounded tuning remains best among the eight benchmark tunings on out-of-sample steps of $-5$, $-1$, $+1$ and $+2$~K and ranks third at $+5$~K.
    \item \textbf{Heat-removal demand explains part of the tracking-error variance and of the valve activity, within wide uncertainty.} Under simultaneous uncertainty the IAE variance of every tuning is attributed mainly to feed concentration and heat-transfer capacity, with bootstrap intervals of $\pm 0.11$ to $\pm 0.17$ that do not separate the two nor the tunings. That attribution is partly carried by the $11.4\%$ of samples in which the setpoint is not attainable: the first-order indices recomputed on the feasible region (the only indices that survive a truncated, dependent box) add up to $0.27$--$0.39$, and the Shapley effects show $UA$ acting mostly through interaction ($0.36$--$0.42$ against a first-order $0.03$--$0.07$). The steady coolant flow the perturbed plant requires explains $18\%$ (CC) to $66\%$ (SIMC) of the rank variance of IAE ($\rho_s^2$) and, for the valve activity, $41\%$ and $48\%$ for AMIGO and SIMC ($\rho_s = +0.64$, $+0.69$) but not for ZN and CC ($\rho_s = -0.18$, $-0.37$).
\end{enumerate}
The released code reproduces every table and figure and accepts new controllers behind the same interface, so that future tuning methods, and nonlinear or learning-based controllers, can be compared with these robust PID baselines under identical plant, actuator and protocol. The generality of the findings is supported by two operating points of one reactor model and one valve model; future work will repeat the ablation, with re-optimized tunings and a Sobol' analysis, on a second reverse-acting process, add valve stiction, and extend the analysis to gain-scheduled PID and to operating points inside the open-loop unstable region.
